\section*{الفصل \ref{ch:b1:crystals-ionic-covalent} --- البلورات (2): الأجسام الصلبة الأيونية والتساهمية والجزيئية}

\begin{solution}{exo:b1:crystals-ionic-covalent:1}
الأنيونات: $8 \times \frac18 + 6 \times \frac12 = 4$؛ والكاتيونات: $12 \times
\frac14 + 1 = 4$. لكل أيون 6 جيران من النوع الآخر (6:6). أربعة
كاتيونات وأربعة أنيونات: \ce{NaCl}.
\end{solution}

\begin{solution}{exo:b1:crystals-ionic-covalent:2}
نصف القطر الرئيسي: $a\sqrt3/2 = 412.3 \times 0.866 = \qty{357.1}{pm}$.
أيون \ce{Cs+} واحد و$8 \times \frac18 = 1$ أيون \ce{Cl-}: وحدة صيغة واحدة.
\end{solution}

\begin{solution}{exo:b1:crystals-ionic-covalent:3}
$x = 102/181 = 0.56$، بين 0.414 و0.732: تناسق ثماني الأوجه
(6)، أي \omterm{def:b1:crystals-ionic-covalent:structure-type}{نمط الملح الصخري}.
\end{solution}

\begin{solution}{exo:b1:crystals-ionic-covalent:4}
فلزي: النحاس. أيوني: كلوريد الصوديوم، وفلوريد الكالسيوم. تساهمي:
الألماس، والكوارتز. جزيئي: ثنائي اليود، والجليد.
\end{solution}

\begin{solution}{exo:b1:crystals-ionic-covalent:5}
على امتداد الحرف، أنيون--كاتيون--أنيون متماسة: $a = 2(r_+ + r_-)$. وعلى
قطر الوجه يقع أنيونان على المسافة $a\sqrt2/2$؛ ويجب ألا
يتداخلا: $a\sqrt2/2 \ge 2r_-$، فيكون $(r_+ + r_-)\sqrt2 \ge 2r_-$، أي
$r_+/r_- \ge \sqrt2 - 1$.
\end{solution}

\begin{solution}{exo:b1:crystals-ionic-covalent:6}
$\rho = 4 \times 58.44/(6.022 \times 10^{23} \times (5.641 \times
10^{-8})^3) = \qty{2.16}{g/cm^3}$، مقابل \qty{2.17}{g/cm^3} مقيسة.
\end{solution}

\begin{solution}{exo:b1:crystals-ionic-covalent:7}
$\rho = 168.36/(6.022 \times 10^{23} \times (4.123 \times 10^{-8})^3) =
\qty{3.99}{g/cm^3}$، وهي القيمة المقيسة.
\end{solution}

\begin{solution}{exo:b1:crystals-ionic-covalent:8}
\ce{Zn-S} $= a\sqrt3/4 = \qty{234.2}{pm}$؛
$\rho = 4 \times 97.44/(6.022 \times 10^{23} \times (5.409 \times
10^{-8})^3) = \qty{4.09}{g/cm^3}$، قريبة من القيمة المقيسة
\qty{4.04}{g/cm^3} للسفاليريت الطبيعي.
\end{solution}

\begin{solution}{exo:b1:crystals-ionic-covalent:9}
\ce{C-C} $= 245.6/\sqrt3 = \qty{141.8}{pm}$؛ وبين الطبقات $669.6/2 =
\qty{334.8}{pm}$. حجم الخلية $\frac{\sqrt3}{2}a^2c = \qty{3.498e-23}{cm^3}$،
$\rho = 4 \times 12.01/(6.022 \times 10^{23} \times 3.498 \times 10^{-23})
= \qty{2.28}{g/cm^3}$، مقابل \qty{3.52}{g/cm^3} للألماس: فالطبقات
متباعدة، لا تمسكها إلا قوى لندن.
\end{solution}

\begin{solution}{exo:b1:crystals-ionic-covalent:10}
انظر برهان \cref{prop:b1:crystals-ionic-covalent:diamond}:
$C = \pi\sqrt3/16 = 0.34$. فالصلادة تأتي من قوة \omterm{def:b1:lewis-resonance-vsepr:lewis}{الروابط التساهمية}
وصلابتها الثلاثية الأبعاد، لا من الرصّ: فلكل
ذرة أربعة جيران فقط، بزوايا ثابتة، وهذا يترك فراغًا
كبيرًا.
\end{solution}

\begin{solution}{exo:b1:crystals-ionic-covalent:11}
\ce{Si-Si} $= a\sqrt3/4 = \qty{235.2}{pm}$، \omterm{def:b1:periodicity:radius}{ونصف القطر التساهمي}
\qty{117.6}{pm}؛ $\rho = 8 \times 28.09/(6.022 \times 10^{23} \times
(5.431 \times 10^{-8})^3) = \qty{2.33}{g/cm^3}$، وهي القيمة المقيسة.
\end{solution}

\begin{solution}{exo:b1:crystals-ionic-covalent:12}
$x = 152/181 = 0.84 > 0.732$: تتنبأ القاعدة \omterm{def:b1:crystals-ionic-covalent:structure-type}{بنمط كلوريد السيزيوم}.
وفي بنية الملح الصخري، \ce{Rb-Cl} $= a/2 = \qty{329.1}{pm}$، قريبة من
$152 + 181 = \qty{333}{pm}$: فالبنية المرصودة متسقة. أما
القاعدة فتعامل الأيونات كرات صلبة وتهمل قابليتها للاستقطاب
والفرق الصغير في الطاقة بين البنيتين (فكلوريد الروبيديوم
يتخذ \omterm{def:b1:crystals-ionic-covalent:structure-type}{نمط كلوريد السيزيوم} تحت الضغط)؛ إنها دليل،
لا قانون.
\end{solution}

\begin{solution}{pb:b1:crystals-ionic-covalent:1}
\textbf{1.} على الرؤوس الثمانية ومراكز الأوجه الستة: $8 \times \frac18 + 6
\times \frac12 = 4$ أيونات.
\textbf{2.} 8 \omterm{def:b1:crystals-metals:site}{مواقع رباعية الأوجه}، كلها مشغولة: 8 أيونات \ce{F-} مقابل 4
أيونات \ce{Ca^{2+}}، بنسبة 1:2، أي \ce{CaF2}.
\textbf{3.} \ce{F-}: 4 أيونات \ce{Ca^{2+}} على رؤوس رباعي أوجه.
\ce{Ca^{2+}}: 8 أيونات \ce{F-} على رؤوس مكعب.
\textbf{4.} $a\sqrt3/4 = 546.3 \times 0.4330 = \qty{236.6}{pm}$.
\textbf{5.} $112 + 131 = \qty{243}{pm}$: في حدود 3\,\%.
\textbf{6.} تكوّن أيونات \ce{F-} ترتيبًا مكعبًا بسيطًا طول حرفه $a/2 =
\qty{273.2}{pm}$، أكبر قليلًا من $2 \times 131 = \qty{262}{pm}$: فهي
تكاد تتماس.
\textbf{7.} $x = 112/131 = 0.85$.
\textbf{8.} $x > 0.732$: تناسق 8.
\textbf{9.} في \omterm{def:b1:crystals-ionic-covalent:structure-type}{نمط كلوريد السيزيوم} عدد الكاتيونات يساوي عدد الأنيونات؛ ومع
أنيونات ضعف الكاتيونات، لا يستطيع إلا نصف مكعبات الأنيونات أن يحوي كاتيونًا.
\textbf{10.} أيونات \ce{F-} الثمانية في الخلية تكوّن مكعبًا من 8 مكعبات صغيرة طول حرفها
$a/2$؛ وتشغل أيونات \ce{Ca^{2+}} مراكز 4 منها، بالتناوب،
كالمربعات السوداء في رقعة شطرنج ثلاثية الأبعاد.
\textbf{11.} $C = \frac43\pi(4 \times 112^3 + 8 \times 131^3)/546.3^3 = 0.61$.
\textbf{12.} \ce{O^{2-}} على \omterm{def:b1:crystals-metals:lattice}{الشبكة} الممركزة الأوجه (الرؤوس ومراكز الأوجه)،
و\ce{Li+} في \omterm{def:b1:crystals-metals:site}{المواقع الرباعية الأوجه} الثمانية.
\textbf{13.} $a\sqrt3/4 = \qty{200.0}{pm}$؛ $59 + 142 = \qty{201}{pm}$.
\textbf{14.} $\rho = 4 \times 29.88/(6.022 \times 10^{23} \times (4.619
\times 10^{-8})^3) = \qty{2.01}{g/cm^3}$، وهي القيمة المقيسة.
\textbf{15.} بلندة الزنك: الألماس هو بلندة الزنك مع الكربون في نوعي
المواضع كليهما.
\textbf{16.} $a\sqrt3/4 = \qty{154.5}{pm}$؛ $\rho = 8 \times
12.01/(6.022 \times 10^{23} \times (3.567 \times 10^{-8})^3) =
\qty{3.52}{g/cm^3}$.
\textbf{17.} في الألماس لا يُملأ إلا نصف \omterm{def:b1:crystals-metals:site}{المواقع الرباعية الأوجه}، وبذرات
بحجم ذرات \omterm{def:b1:crystals-metals:lattice}{الشبكة}، لا تتماس إلا على امتداد الروابط:
$C = 0.34$. أما في الفلوريت فكل المواقع مملوءة والكاتيونات الصغيرة
تملأ الفجوات بين الأنيونات الكبيرة: $C = 0.61$.
\textbf{18.} $40.08 + 2 \times 19.00 = \qty{78.08}{g/mol}$.
\textbf{19.} $4 \times 78.08/(6.022 \times 10^{23}) = \qty{5.186e-22}{g}$.
\textbf{20.} $(5.463 \times 10^{-8})^3 = \qty{1.630e-22}{cm^3}$.
\textbf{21.} غياب أيون \ce{F-} واحد في كل مئة خلية ينقص الكتلة بمقدار $19.00/(100
\times 312.3)$، أي 0.06\,\%: لا يظهر عند ثلاثة أرقام (ويجب أن يغيب
أيون \ce{Ca^{2+}} أيضًا، وإلا لم تكن \omterm{def:b1:crystals-metals:crystal}{البلورة} متعادلة).
\textbf{22.} $\rho = 5.186 \times 10^{-22}/1.630 \times 10^{-22} =
\textbf{\qty{3.18}{g/cm^3}}$، وهي الكتلة الحجمية المقيسة تمامًا: فنموذج
الخلية، المبني على قياس واحد لوسيط $a$ بالأشعة السينية، يفسّر كتلة
المعدن.
\end{solution}
